\documentclass[11pt,letterpaper]{article}

\usepackage[margin=1in]{geometry}
\usepackage[T1]{fontenc}
\usepackage[utf8]{inputenc}
\usepackage{lmodern}
\usepackage{fancyhdr}

\setlength{\parindent}{0pt}
\setlength{\parskip}{0.65em}
\newcommand{\checkmarkbox}{\fbox{\rule{0pt}{0.9ex}\rule{0.9ex}{0pt}}}

\pagestyle{fancy}
\fancyhf{}
\fancyhead[L]{\small Deep Learning and Reinforcement Learning}
\fancyhead[R]{\small Final Project}
\fancyfoot[C]{\small\thepage}
\renewcommand{\headrulewidth}{0.3pt}
\renewcommand{\footrulewidth}{0pt}

\begin{document}

\begin{titlepage}
  \thispagestyle{empty}
  \vspace*{1.25in}
  \rule{\textwidth}{2pt}
  \vspace{0.35in}

  {\Huge\bfseries Final Project Report\par}
  \vspace{0.2in}
  {\Large Deep Learning and Reinforcement Learning\par}

  \vspace{0.55in}
  \begin{minipage}{0.86\textwidth}
    \large
    This document presents the objectives, report structure, deliverables,
    and evaluation criteria for the final project. It is intended to guide
    the preparation of a concise, stakeholder-ready analytical report.
  \end{minipage}

  \vfill
  \begin{tabular}{@{}ll@{}}
    \textbf{Deliverable:} & Written analytical report submitted as a PDF\\
    \textbf{Audience:} & Senior business and analytics stakeholders\\
    \textbf{Focus:} & Deep learning or reinforcement learning application
  \end{tabular}

  \vspace{0.55in}
  \rule{\textwidth}{2pt}
\end{titlepage}

\tableofcontents
\newpage

\section{Project Overview}

One of the main objectives of this course is to provide hands-on experience
in communicating insightful and impactful findings to stakeholders. In this
project, you will use the tools and techniques learned throughout the course
to apply deep learning to a task of your choosing.

The application may be supervised or unsupervised and may involve
classification, image analysis, text analysis, or another suitable problem.
You will refine your model, compare alternatives, and communicate the
insights found during model development. The final-project questions will be
derived from the selected case study. You may submit the assignment using
either the AI-graded or peer-reviewed option.

\section{Choosing a Data Set}

Before beginning the analysis, select a data set that is meaningful to you.
Consider a public data set similar to information available in your workplace,
a data set you have always wanted to analyze, sports data, data from a
datathon, or a data-for-good project.

Sample data sets used in the course's Python notebook demonstrations may also
be used. You may optionally use the discussion board to brainstorm with peers
and identify suitable data sets. Participation in the discussion board is not
required.

\section{Final Project Report Requirements}

The primary report should emphasize analytical judgment, a transparent
thought process, and actionable next steps. Write for a senior audience, such
as a Chief Data Officer or Head of Analytics. The report must include the
following sections.

\begin{enumerate}
  \item \textbf{Objective and business value.} State the main objective of
    the analysis, identify whether the model uses a specific deep learning or
    reinforcement learning algorithm, and explain the benefits for the
    business or relevant stakeholders.

  \item \textbf{Data description.} Describe the selected data set, summarize
    its attributes, and explain what the analysis is intended to accomplish.

  \item \textbf{Exploration and preparation.} Summarize the data exploration
    process and document actions taken for data cleaning, transformation, or
    feature engineering.

  \item \textbf{Model variations.} Summarize training for at least three
    variations of the selected deep learning model. Variations may use
    different architectures, clustering techniques, hyperparameters, or other
    justified choices.

  \item \textbf{Final model recommendation.} Explain which model you
    recommend and why it best fits the stated needs, considering accuracy,
    interpretability, operational constraints, and explainability.

  \item \textbf{Key findings and insights.} Walk the reader through the main
    findings from the modeling exercise and connect them explicitly to the
    project objectives.

  \item \textbf{Next steps.} Suggest practical improvements, such as
    revisiting the model, collecting additional observations, or adding
    features that could improve performance or usefulness.
\end{enumerate}

\section{Deliverable and Submission Guidelines}

Submit a written analytical report as a PDF. The report should focus on
insights, findings, and recommendations rather than presenting an undigested
record of experimentation. Include visualizations when they support the
analysis and help the reader understand the data, model behavior, or
recommendations.

\subsection{Recommended Report Structure}

\begin{itemize}
  \item Executive summary
  \item Business objective and stakeholder value
  \item Data description and preparation
  \item Exploratory analysis
  \item Model design and comparison
  \item Recommended model
  \item Key findings and limitations
  \item Recommendations and next steps
\end{itemize}

\section{Evaluation Criteria}

The project will be evaluated using the following criteria:

\begin{itemize}
  \item A clear and comprehensive description of the data
  \item A well-defined statement of the primary analysis objective(s)
  \item A comparison of multiple deep learning models with a justified final
    model selection
  \item Clear presentation of key findings aligned with the objective(s)
  \item Identification of model limitations and a well-articulated plan for
    future improvement
\end{itemize}

\section{Submission Checklist}

\begin{table}[htbp]
  \centering
  \caption{Final review checklist}
  \label{tab:checklist}
  \begin{tabular}{@{}p{0.68\textwidth}p{0.18\textwidth}@{}}
    \hline
    \textbf{Requirement} & \textbf{Complete}\\
    \hline
    Objective and stakeholder value are clearly stated & \checkmarkbox\\
    Data set, attributes, and intended outcome are described & \checkmarkbox\\
    Exploration, cleaning, and feature engineering are documented & \checkmarkbox\\
    At least three model variations are compared & \checkmarkbox\\
    Final model is recommended with supporting justification & \checkmarkbox\\
    Key findings are connected to the project objective & \checkmarkbox\\
    Limitations and practical next steps are included & \checkmarkbox\\
    Report is submitted as a polished PDF & \checkmarkbox\\
    \hline
  \end{tabular}
\end{table}

\end{document}
